\section{方式四：AR标记识别与队列跟随（第10组）}
\subsection{实验目标与总体方案}
本组使用三台TurtleBot2纵列跟随：一号车以点云视觉跟踪TonyPi，二号车跟随一号车，三号车跟随二号车。二、三号车通过Kinect识别前车背部AR标签，根据相对位置控制底盘；深度图转换得到的扫描用于近障响应。该方案不使用预建地图和全局路径规划。

二、三号车使用\code{team_turtlebot_follow_plus}中的C++节点。它们直接根据标签和扫描产生速度，区别于通过前车AMCL位姿和地图规划进行跟随的方案。

一号车和二号车尾部分别粘贴ID 0和ID 1标签，三号车尾部未贴标签。二号车识别一号车的ID 0标签，三号车识别二号车的ID 1标签。一号车通过点云数据识别并跟踪TonyPi，不使用AR标签。以下功能包和启动参数说明适用于二、三号车。

\subsection{感知输入与目标定位}
功能包\code{team\_turtlebot\_follow\_plus}使用ROS Kinetic/C++实现，位于\code{turtlebot/follow4/team\_turtlebot\_follow\_plus/}。\code{ar\_track\_alvar}根据Kinect RGB图像和相机内参发布标签位姿；\code{depthimage\_to\_laserscan}将深度图转换为\code{/scan}。跟随节点订阅\code{/ar\_pose\_marker}与\code{/scan}，向\code{/mobile\_base/commands/velocity}发布速度指令。


相机输入、内参和标签尺寸应成套核对。二、三号车分别只选择ID 0和ID 1；其他标签即使进入画面，也不应改变直接前车的身份。深度扫描负责近障判断，标签位姿负责跟随方向，两种观测承担不同任务。

\subsection{运动决策与控制方法}
设相机光学坐标系中的标签位置为$(x,z)$，其中$x$向图像右侧为正、$z$指向相机前方，则标签距离与方位角为
\begin{equation}
 r=\sqrt{x^2+z^2},\qquad \theta=\operatorname{atan2}(x,z).
\end{equation}
控制节点以10\,Hz运行，按距离误差比例调节前进速度，按方位角调节转向速度；线速度限制在零至设定上限之间，目标距离不足时不倒车，标签偏离车头方向较大时暂停前进。

避障逻辑将扫描数据划分为前方及左右侧扇区。前方出现近障时，节点选择有足够扫描净空的一侧并保持该转向选择；距离达到停车阈值或所选侧净空不足时停止前进。扫描数据超时则发布零速度。目标首次出现前车辆保持停止；已识别目标丢失后，车辆低速原地搜索，超过设定时限仍未重获目标则停止。该方法依赖标签保持在相机视野内，不能沿被遮挡路线推断前车位置。

\subsection{运行流程与参数设置}
二、三号车使用相同的跟随和避障参数，目标标签ID分别设为0和1。表\ref{tab:f4-param-meaning}列出实际启动值及参数用途。

{\small\renewcommand{\arraystretch}{1.05}\setlength{\tabcolsep}{3pt}
\begin{longtable}{p{31mm}p{10mm}p{27mm}p{75mm}}
\caption{二、三号车运行参数及含义}\label{tab:f4-param-meaning}\\
\toprule
参数 & 单位 & 二号车 / 三号车 & 含义\\\midrule
\endfirsthead
\caption[]{二、三号车运行参数及含义（续）}\\
\toprule
参数 & 单位 & 二号车 / 三号车 & 含义\\\midrule
\endhead
\code{marker\_id} & -- & 0 / 1 & 目标标签编号。节点只跟随该ID对应的标签。\\
\code{marker\_size} & cm & 14.0 / 14.0 & 标签黑色方框的实际边长，用于AR距离估计，应与打印尺寸一致。\\
\code{desired\_distance} & m & 0.40 / 0.40 & Kinect到标签的目标距离，不等同于两车外壳间距。程序不倒车；标签近于设定距离时停止前进。\\
\code{max\_linear\_speed} & m/s & 0.35 / 0.35 & 最大前进速度。实际速度随距离误差及避障控制变化。初次低速测试设为0.05--0.10\,m/s。\\
\code{max\_angular\_speed} & rad/s & 0.25 / 0.25 & 正常跟随及近障转向的角速度上限。\\
\code{avoid\_distance} & m & 0.38 / 0.38 & 前方障碍进入该距离后，开始减速并向扫描较空的一侧偏转。\\
\code{stop\_distance} & m & 0.20 / 0.20 & 前方障碍达到该距离后，线速度置零；侧向净空足够时向较空的一侧转向。\\
\code{side\_clearance} & m & 0.40 / 0.40 & 转向侧所需的最小扫描净空。两侧均不足时车辆停止。\\
\code{avoid\_release\_margin} & m & 0.20 / 0.20 & 障碍距离超过\code{avoid\_distance}加该余量后，退出避障状态并恢复跟随。\\
\code{obstacle\_turn\_speed} & rad/s & 0.20 / 0.20 & 达到停车距离后，车辆向选定空旷侧转向时使用的角速度上限。\\
\code{marker\_timeout} & s & 0.8 / 0.8 & 标签超过该时长没有新检测结果后，节点开始目标搜索。\\
\code{lost\_search\_timeout} & s & 16 / 16 & 原地搜索目标的最长时间；超时仍未重获目标时停止。\\
\code{search\_angular\_speed} & rad/s & 0.20 / 0.20 & 目标丢失后原地搜索时的角速度，且不超过\code{max\_angular\_speed}。\\
\bottomrule
\end{longtable}}


运行时先编译并加载工作空间，再分终端启动底盘、Kinect和深度转扫描节点；确认扫描连续更新后，分别以ID 0和ID 1启动二、三号车。初次测试先以0.05--0.10\,m/s观察响应，再结合场地调整；表中的0.35\,m/s是启动配置的速度上限。完整命令见附录\ref{app:queue-commands}，便于调试时查阅。

\subsection{实验现象与结果分析}
三车联调结果见表\ref{tab:f4-results}。本次完成了直行、转弯和依次停车，没有专门布置障碍物，也没有出现目标丢失，因此避障与重捕获还需要增加相应场景测试。

\noindent\begin{minipage}{\linewidth}\centering\small
\captionof{table}{方式四实验记录}\label{tab:f4-results}
\vspace{-3pt}\renewcommand{\arraystretch}{1.05}\setlength{\tabcolsep}{4pt}
\begin{tabularx}{\linewidth}{p{39mm}Yp{22mm}}\toprule
测试项目 & 实际情况 & 结果\\\midrule
直行跟随 & 队列间距保持稳定，速度适中，车辆间未碰撞 & 完成\\
转弯跟随 & 采用圆弧轨迹转弯，保持后车对前车标签的观测 & 完成\\
障碍响应 & 最终场地未专门布置障碍；未碰撞围栏，但围栏及长时间停留的前车曾触发误避障 & 有待优化\\
目标保持 & 本次联调未发生目标丢失，重捕获逻辑未实际触发 & 未覆盖\\
终点停车 & 三车依次到达终点并平稳停靠 & 完成\\
\bottomrule\end{tabularx}
\end{minipage}

赛道围栏未发生碰撞，但当前车低速行驶时，后车曾把围栏识别为障碍并触发避障；当前车长时间停留在近处时，后车也曾将其判断为障碍，继而原地转向。后续需结合赛道宽度和前后车间距调整避障距离及状态切换条件，并在围栏邻近、前车慢行和停车场景下复测。测试期间目标未丢失，目标重捕获功能尚无实际触发记录。

\subsection{小组协作与改进总结}
徐杰鸿负责跟随程序与参数，王振宇负责标签和车辆配置，林琢严配合控制优化与场地布置，何爽负责障碍场景、运行记录和报告。我们共同完成三车联调，遇到围栏和停车前车引起的误避障时，一起观察扫描和车辆位置再调整参数。详细分工见第\ref{sec:division}节。




本组以点云视觉跟踪和AR码标签跟随组成三车队列：一号车跟踪TonyPi，二号车识别ID 0，三号车识别ID 1；标签边长为14\,cm。联调中三车完成直行、圆弧转弯和依次停车，车辆未发生碰撞或目标丢失。围栏和慢行前车引发的误避障仍需通过参数调整和复测处理；本次未进行专门障碍布置，不能据此判定复杂障碍下的避障性能。
